
\documentclass[a4paper,11pt]{article}
\usepackage[a4paper, margin=8em]{geometry}

% usa i pacchetti per la scrittura in italiano
\usepackage[french,italian]{babel}
\usepackage[T1]{fontenc}
\usepackage[utf8]{inputenc}
\frenchspacing 

% usa i pacchetti per la formattazione matematica
\usepackage{amsmath, amssymb, amsthm, amsfonts}

% usa altri pacchetti
\usepackage{gensymb}
\usepackage{hyperref}
\usepackage{standalone}

% imposta il titolo
\title{Appunti Comunicazioni Numeriche}
\author{Luca Seggiani}
\date{2026}

% imposta lo stile
% usa helvetica
\usepackage[scaled]{helvet}
% usa palatino
\usepackage{palatino}
% usa un font monospazio guardabile
\usepackage{lmodern}

\renewcommand{\rmdefault}{ppl}
\renewcommand{\sfdefault}{phv}
\renewcommand{\ttdefault}{lmtt}

% disponi il titolo
\makeatletter
\renewcommand{\maketitle} {
	\begin{center} 
		\begin{minipage}[t]{.8\textwidth}
			\textsf{\huge\bfseries \@title} 
		\end{minipage}%
		\begin{minipage}[t]{.2\textwidth}
			\raggedleft \vspace{-1.65em}
			\textsf{\small \@author} \vfill
			\textsf{\small \@date}
		\end{minipage}
		\par
	\end{center}

	\thispagestyle{empty}
	\pagestyle{fancy}
}
\makeatother

% disponi teoremi
\usepackage{tcolorbox}
\newtcolorbox[auto counter, number within=section]{theorem}[2][]{%
	colback=blue!10, 
	colframe=blue!40!black, 
	sharp corners=northwest,
	fonttitle=\sffamily\bfseries, 
	title=~\thetcbcounter: #2, 
	#1
}

% disponi definizioni
\newtcolorbox[auto counter, number within=section]{definition}[2][]{%
	colback=red!10,
	colframe=red!40!black,
	sharp corners=northwest,
	fonttitle=\sffamily\bfseries,
	title=~\thetcbcounter: #2,
	#1
}

% U.D.M
\newcommand{\amp}{\ensuremath{\mathrm{A}}}
\newcommand{\volt}{\ensuremath{\mathrm{V}}}
\newcommand{\meter}{\ensuremath{\mathrm{m}}}
\newcommand{\second}{\ensuremath{\mathrm{s}}}
\newcommand{\farad}{\ensuremath{\mathrm{F}}}
\newcommand{\henry}{\ensuremath{\mathrm{H}}}
\newcommand{\siemens}{\ensuremath{\mathrm{S}}}

% circuiti
\usepackage{circuitikz}
\usetikzlibrary{babel}

% disegni
\usepackage{pgfplots}
\pgfplotsset{width=10cm,compat=1.9}

% disponi codice
\usepackage{listings}
\usepackage[table]{xcolor}

\lstdefinestyle{codestyle}{
		backgroundcolor=\color{black!5}, 
		commentstyle=\color{codegreen},
		keywordstyle=\bfseries\color{magenta},
		numberstyle=\sffamily\tiny\color{black!60},
		stringstyle=\color{green!50!black},
		basicstyle=\ttfamily\footnotesize,
		breakatwhitespace=false,         
		breaklines=true,                 
		captionpos=b,                    
		keepspaces=true,                 
		numbers=left,                    
		numbersep=5pt,                  
		showspaces=false,                
		showstringspaces=false,
		showtabs=false,                  
		tabsize=2
}

\lstdefinestyle{shellstyle}{
		backgroundcolor=\color{black!5}, 
		basicstyle=\ttfamily\footnotesize\color{black}, 
		commentstyle=\color{black}, 
		keywordstyle=\color{black},
		numberstyle=\color{black!5},
		stringstyle=\color{black}, 
		showspaces=false,
		showstringspaces=false, 
		showtabs=false, 
		tabsize=2, 
		numbers=none, 
		breaklines=true
}

\lstdefinelanguage{javascript}{
	keywords={typeof, new, true, false, catch, function, return, null, catch, switch, var, if, in, while, do, else, case, break},
	keywordstyle=\color{blue}\bfseries,
	ndkeywords={class, export, boolean, throw, implements, import, this},
	ndkeywordstyle=\color{darkgray}\bfseries,
	identifierstyle=\color{black},
	sensitive=false,
	comment=[l]{//},
	morecomment=[s]{/*}{*/},
	commentstyle=\color{purple}\ttfamily,
	stringstyle=\color{red}\ttfamily,
	morestring=[b]',
	morestring=[b]"
}

% disponi sezioni
\usepackage{titlesec}

\titleformat{\section}
	{\sffamily\Large\bfseries} 
	{\thesection}{1em}{} 
\titleformat{\subsection}
	{\sffamily\large\bfseries}   
	{\thesubsection}{1em}{} 
\titleformat{\subsubsection}
	{\sffamily\normalsize\bfseries} 
	{\thesubsubsection}{1em}{}

% disponi alberi
\usepackage{forest}

\forestset{
	rectstyle/.style={
		for tree={rectangle,draw,font=\large\sffamily}
	},
	roundstyle/.style={
		for tree={circle,draw,font=\large}
	}
}

% disponi algoritmi
\usepackage{algorithm}
\usepackage{algorithmic}
\makeatletter
\renewcommand{\ALG@name}{Algoritmo}
\makeatother

% disponi numeri di pagina
\usepackage{fancyhdr}
\fancyhf{} 
\fancyfoot[L]{\sffamily{\thepage}}

\makeatletter
\fancyhead[L]{\raisebox{1ex}[0pt][0pt]{\sffamily{\@title \ \@date}}} 
\fancyhead[R]{\raisebox{1ex}[0pt][0pt]{\sffamily{\@author}}}
\makeatother

\begin{document}
% sezione (data)
\section{Lezione del 23-04-26}

% stili pagina
\thispagestyle{empty}
\pagestyle{fancy}

% testo
Alla fine della lezione 16 avevamo visto i sistemi di \textit{più} variabili aleatorie, e la definizione di \textit{probabilità congiunta}.

\subsection{Correlazione e covarianza}
Il comportamento statistico di una variabile aleatoria $X$ può essere caratterizzato in maniera incompleta ma talvolta sufficiente da alcuni parametri caratteristici (che avevamo visto in 13.1), come valor medio e varianza.

Anche per i sistemi di più variabili aleatorie, ad esempio $X$ e $Y$, possiamo determinare alcuni parametri statistici semplificati che forniscono utili indicazioni per la compresnsione del loro comportamento statistico congiunto. Queste sono:
\begin{itemize}
	\item La \textbf{correlazione} $r_{XY}$:
$$
r_{XY} = E \{ XY \} = \int_{-\infty}^{\infty} \int_{-\infty}^{\infty} x y f_{XY} (x, y) \, dx dy
$$

Se $r_{XY} = 0$, si dice che le due variabili sono \textbf{ortogonali};

\item La \textbf{covarianza} $c_{XY}$:
$$
	c_{XY} = E\{ (X - \eta_X) (Y - \eta_Y) \} = \int_{-\infty}^{\infty} \int_{-\infty}^{\infty} (x - \eta_X) (y - \eta_Y) f_{XY} (x, y) \, dx dy = r_{XY} - \eta_X \eta_Y
$$

Se $c_{XY} = 0$, si dice che le due variabili sono \textbf{incorrelate}.
La covarianza $c_{XY}$ è un parametro statistico molto importante, in quanto accerta se tra le due variabili $X$ e $Y$ esiste una relazione di dipendenza di tipo lineare. In particolare, misura la tendenza di variazione congiunta (covarianza) delle due. Se la covarianza è grande è positiva, le due variabili $X$ e $Y$ si discostano dal valor medio nella stessa direzione, cioè le quantità $X - \eta_X$ e $Y - \eta_Y$ tendono ad avere lo stesso segno.
\end{itemize}

Da correlazione e covarianza ricaviamo un terzo indice, che ha in maniera più distillata lo stesso significato della covarianza: il \textbf{coefficiente di correlazione} (o \textit{covarianza normalizzata}):
$$
\rho_{XY} = E \left\{ \frac{X - \eta_X}{\sigma_X} \cdot \frac{Y - \eta_Y}{\sigma_Y} \right\} = \frac{c_{XY}}{\sigma_X \sigma_Y} = \frac{r_{XY} - \eta_X \eta_Y}{\sigma_X \sigma_Y}
$$

Vediamo le proprietà di questa definizione:
\begin{enumerate}
	\item Poichè $c_{XY}^2 \leq \sigma_X^2 \sigma_Y^2$, si ha $|\rho_{XY}| < 1$;
	\item Se le variabili sono incorrelate, $\rho_{XY} = 0$;
	\item Se $Y = aX + b$ (con $a \neq 0$), allora chiaramente $|\rho_{XY}| = 1$.
\end{enumerate}
Con riferimento alla terza proprietà elencata, si può dimostrare che se viceversa il coefficiente di correlazione è 1, le variabili aleatorie sono linearmente dipendenti. Per $|\rho_{XY}| = 1$ è possibile fare una previsione perfetta, cioè noto $X$ calcolare $Y$ come $Y = aX + b$.

\par\smallskip

Notiamo quindi come la varianza ci permette di calcolare velocemente la varianza di due variabili aleatorie:
$$
\mathrm{Var}(X + Y) = \mathrm{Var}(x) + \mathrm{Var}(Y) + c_{XY}
$$
in quanto:
$$
\mathrm{Var}(X + y) = E\left[(X + Y)^2\right] - E^2 \left[ X + Y \right]
$$
che per linearità:
$$
= E\left[X^2\right] + E\left[Y^2\right] + 2 E\left[XY\right] - E^2\left[X\right] - E^2\left[Y\right] - 2E\left[X\right]E\left[Y\right]
$$
$$
= E\left[X^2\right] - E^2\left[X\right] + E\left[Y^2\right] - E^2\left[Y\right] + 2 \left( E \left[ XY \right] - E\left[X\right] E\left[Y\right] \right)
$$
che è la tesi.

\subsubsection{Correlazione e covarianza di variabili indipendenti}
Notiamo cosa succede a correlazione e covarianza nel caso di variabili aleatorie indipendenti. Iniziamo dalla correlazione:
$$
r_{XY} = E \left\{XY\right\}
= \int_{-\infty}^{\infty} \int_{-\infty}^{\infty} x y f_{XY} (x, y) \, dx dy 
= \int_{-\infty}^{\infty} \int_{-\infty}^{\infty} x y f_{X} f_{Y} (x, y) \, dx dy
$$
$$
= \int_{-\infty}^{\infty} x f_X(x) \, dx + \int_{-\infty}^{\infty} y f_Y(y) \, dy = \eta_X \eta_Y
$$
Quindi si ha semplicemente il prodotto dei valor medi. Da questo si ricavano immediatamente covarianza e coefficiente di correlazione come nulli:
$$
c_{XY} = r_{XY} - \eta_X \eta_Y = 0 , \quad \rho_{XY} = \frac{c_{XY}}{\sigma_X \sigma_Y} = 0
$$

Abbiamo quindi che a variabili aleatorie indipendenti corrisponde che le due variabili sono incorrelate. Cioè:
$$
X, Y \text{ indipendenti} \implies X, y \text{ incorrelate}
$$
Notiamo che non vale assolutamente il contrario, cioè non vale:
$$
X, Y \text{ incorrelate} \mathrel{{\ooalign{\hidewidth$\not\phantom{=}$\hidewidth\cr$\implies$}}} X, y \text{ indipendenti}
$$

\subsection{Sistemi di variabili congiuntamente gaussiane}
Chiamiamo due variabili aleatorie continue $X$ e $Y$ \textbf{conginutamente gaussiane} se la loro funzione di densità di probabilità congiunta ha la seguente espressione:
$$
f_{XY} (x, y) = \frac{1}{2 \pi \sigma_X \sigma_Y \sqrt{1 - \rho^2}} e^{- \frac{1}{2(1 - \rho^2)} Q(x, y)}
$$
dove $Q(x, y)$ è:
$$
Q(x, y) = \frac{(x - \eta_X)^2}{\sigma_X^2} - \frac{2 \rho(x - \eta_X)(y - \eta_Y)}{\sigma_X \sigma_Y} + \frac{(y - \eta_Y)^2}{\sigma_Y^2}
$$
e $\rho$ è il coefficiente di correlazione:
$$
\rho_{XY} = \frac{c_{XY}}{\sigma_X \sigma_Y}
$$

L'idea fondamentale è che se $X$ e $Y$ sono congiuntamente gaussiane, anche le distribuzioni di $X$ e $Y$ prese singolarmente sono gaussiane.

Proprietà delle variabili aleatorie continue congiuntamente gaussiane è che se due variabili aleatorie sono congiuntamente gaussiane e incorrelate, allore si dimostra che sono indipendenti. Infatti, se $\rho = 0$, la distribuzione di probabilità congiuntà diventa il prodotto delle distribuzioni di probabilità marginali, che è la definizione di indipendenza:
$$
f_{XY} (x, y) = \frac{1}{2 \pi \sigma_X \sigma_Y} e^{ -\frac{(x - \eta_X)^2}{2\sigma_X^2} -\frac{(y - \eta_Y)^2}{2\sigma_Y^2} }
= \frac{1}{2 \pi \sigma_X} e^{ -\frac{(x - \eta_X)^2}{2\sigma_X^2} } \frac{1}{2 \pi \sigma_Y} e^{ -\frac{(y - \eta_Y)^2}{2\sigma_Y^2} }
$$

\subsection{Teorema del limite centrale}
Consideriamo una variabile aleatoria $Y_n = \sum_{i = 1}^n X_i$. Supponiamo che le $n$ variabili aleatorie $X_i$ siano indipendenti e con la stessa funzione densità di probabilità con valor medio $\eta$ e varianza $\sigma^2$. Allora sarà vero che:
$$
E \{ Y_n \} = \eta_n = n \eta, \quad E\{ (Y_n - \eta_n)^2 \} = \sigma_n^2 = n \sigma^2
$$

Se consideriamo un $n$ via via crescente, la densità di probabilità della variabile aleatoria normalizzata:
$$
S_n = \frac{Y_n - \eta_n}{\sigma_n} = \frac{Y_n - n \eta}{\sqrt{n} \sigma}
$$
tende ad una densità normale standard, cioè:
$$
\lim_{n \rightarrow \infty} f_{S_n} (s) = f_Z(s) = \frac{1}{\sqrt{2\pi}} e^{-\frac{s^2}{2}}
$$

Questa affermazione è quella del teorema del limite centrale, che afferma che esperimenti ripetuti con una qualsiasi distribuzione di probabilità convergono ad una distribuzione gaussiana con media $n \eta$ e deviazione standard $\sqrt{n} \sigma$.

Più formalmente:
\begin{theorem}{Teorema del limite centrale}
Data una variabile aleatoria $Y_n$ definita come:
$$
Y_n = \sum_{i=1}^n X_i
$$
cioè come iterazione di una variabile aleatoria $X_i$ con valor medio $\eta$ e varianza $\sigma$, presa la variabile standardizzata: 
$$
S_n = \frac{Y_n - n \eta}{\sqrt{n} \sigma}
$$
la sua funzione densità di probabilità tende alla gaussiana standardizzata:
$$
\lim_{n \rightarrow \infty} f_{S_n} (s) = f_Z(s) = \frac{1}{\sqrt{2\pi}} e^{-\frac{s^2}{2}}
$$
\end{theorem}

\subsection{Processi aleatori}
Dato un esperimento casuale di modello di probabilità assegnato, ad ogni suo risultato $\omega_i$ associamo una funzione $\omega_i$ assegnamo una funzione reale $x(t, \omega)$ sulla variabile (solitamente tempo) $t$. Abbiamo così definito un insieme di funzioni $X(t, \omega)$, dette \textbf{processi aleatori}. 

\begin{definition}{Processo aleatorio}
Chiamiamo processo aleatorio una funzione reale:
$$
x(t, \omega)
$$
associata ad ogni risultato $\omega$ di un esperimento casuale.
\end{definition}

Per brevità rappresenteremo i processi aleatori con la sola dipendenza in $t$, cioè $X(t)$.

Riserviamo $X(t, \omega)$ per la \textit{variabile aleatoria} stessa, e $x(t, \omega)$ per la \textit{realizzazione} di questa (ci sarà più chiaro cos'è sotto).
\begin{itemize}
	\item Fissato un determinato istante di tempo $t$, possiamo dire che $X(t, \omega)$ è una \textit{variabile aleatoria} su ogni risultato $\omega_i$ dell'esperimento. Ad ogni $\omega_i$ viene quindi associato il valore numerico della corisspondente realizzazione a tale istante;
	\item Fissato il risultato $\omega_i$, la $x(t, \omega_i)$ è una possibile realizzazione della variabile aleatoria nel tempo.
\end{itemize}

Ciò che stiamo facendo è quindi parametrizzare le nostre variabili aleatorie sulla variabile $t$, solitamente tempo, in maniera tale che campionare la variabile aleatoria non restituisce più un valore costante ma uno che si evolve su al variare di $t$. Notiamo che le singole funzioni $x(t, \omega)$ sono funzioni deterministiche, e la casualità risiede solo nella presentazione di un particolare risultato dell’esperimento. In particolare, la $x(t, \omega)$ che si osserva in una data prova dell’esperimento aleatorio prende il nome di realizzazione del processo.

\begin{itemize}
	\item Se si fissano due istanti distinti $t_1$ e $t_2$, si ottengono due variabili aleatorie distinte $X(t_1)$ e $X(t_2)$, che costituiscono un sistema di due variabili aleatorie, ovvero il \textit{vettore aleatorio} $\mathbf{X} = ( X(t_1) \, X(t_2) )$;
	\item Analogamente, fissando $N$ istanti $t_i$, si ottiene un vettore di $N$ variabili aleatorie $\mathbf{X} = ( X(t_1) \, ... \, X(t_N) )$
\end{itemize}

La descrizione statistica del processo implica perciò la conoscenza della legge di distribuzione di tutti i possibili sistemi così formati.

Riassumendo, $X(t, \omega)$ può declinarsi in:
\begin{itemize}
	\item $X(t, \omega)$ è un insieme di funzioni delle variabili $t$ e $\omega$ (processo aleatorio);
	\item $x(t, \omega_i)$ è una funzione deterministica della variabile $t$, detta \textit{realizzazione} del processo;
	\item $X(t_k, \omega)$ è una \textit{variabile aleatoria} ($t = t_k$ fisso, $\omega$ variabile dipendente dai risultati di un esperimento casuale);
	\item $x(t_k, \omega_i)$ è quindi un numero reale.
\end{itemize}
Vediamo allora come cambiamo fra $X$ maiuscola e $x$ minuscola: solitamente si associa la maiuscola a variabili aleatorie e la minuscola a variabili determinstiche. Vediamo quindi che i n molte applicazioni i risultati dell’esperimento sono già delle forme d’onda: in tal caso non vi è più distinzione tra risultato e realizzazione (l'$\omega_i$ corrisponde direttamente alla $x(t, \omega_i)$ realizzazione).

\subsubsection{Descrizione statistica di un processo aleatorio}
Esistono più modi per descrivere un processo aleatorio:
\begin{enumerate}
\item
Un processo $X(t)$ si dice \textbf{statisticamente determinato} se sono note le sue funzioni di distribuzione:
$$
F_X(x_1, ..., x_N; \, t_1, ..., t_N) = \mathrm{Pr} \left\{ X(t_1) \leq x_1, ..., X(t_N) \leq x_N,  \right\}
$$
dato $N$ e la $N$-upla di istanti temporali $t_i$. Abbiamo quindi che nota la distribuzione di ordine $N$ è possibile ricavare tutte le distribuzione di ordine inferiore mediante le regole marginali (non vale il viceversa).

La funzione di distribuzione di ordine $N$ del processo è la funzione di distribuzione del vettore di variabili aleatorie:
$$
\mathbf{X} = \begin{pmatrix}
	X(t_1) & ... & X(t_N)
\end{pmatrix} 
$$
ottenuto fissando $N$ istanti $t_1, \, ..., \, t_n$.

Il comportamento statistico di un processo stocastico è completamente determinato quando sono note le distribuzioni di tutti i possibili ordini, tuttavia in
alcune applicazioni è sufficiente conoscere alcune statistiche dei primi due ordini, cioè la descrizione in \textbf{potenza} del processo.

\item
Un processo $X(t)$ si dice \textbf{parametrico} quando può essere specificato attraverso la forma delle sue realizzazioni, che dipende parametricamente da un certo numero di variabili aleatorie:
$$
X(t) = s(t, \Omega_1, ..., \Omega_K)
$$
In questo caso la caratterizzazione statistica completa del processo richiede la conoscenza della densità di probabilità congiunta dei parametri aleatori, cioè la:
$$
f_\Omega(\theta_1, \theta_2, ... \theta_k)
$$

\end{enumerate}

\end{document}
